\exercise{11.5}{Subsequential limits of an enumeration of a rational interval}
\goal Let $(q_n)$ enumerate all rational numbers in $(0,1]$.
Find \textup{(a)} its set of subsequential limits and
\textup{(b)} its limit superior and limit inferior.

\idea Every term lies between $0$ and $1$, so a subsequential limit
must lie in $[0,1]$. Conversely, every neighborhood of a point of
$[0,1]$ contains infinitely many distinct rational values in $(0,1]$.
Removing a finite prefix of the enumeration cannot remove all those
values. This will allow a construction with strictly increasing indices.

\step{1. Explain the role of enumeration and distinct values}
The enumeration hypothesis says that each rational number in $(0,1]$
occurs somewhere in the list. It does not require a particular rational
number to occur infinitely often. For example, the value $1$ may occur
only once. The proof below uses infinitely many \emph{different}
rational values near a proposed limit, so it works both with and
without repeated terms.

Every nonempty open interval $(c,d)$ contains infinitely many rational
numbers. To prove this from density, suppose its rational numbers
formed a finite set. Density ensures the set is nonempty, so it would
have a largest member $r<d$. Density applied to $(r,d)$ would produce
a rational number larger than $r$ still belonging to $(c,d)$,
contradicting the choice of $r$. Thus there are infinitely many distinct
rational numbers in the interval.

\step{2. Part (a): Find infinitely many eligible values near every target}
Fix $t\in[0,1]$ and $\varepsilon>0$. Consider
\[
 c=\max\{0,t-\varepsilon\},
 \qquad d=\min\{1,t+\varepsilon\}.
\]
We have $c<d$. At $t=0$, this follows from
$c=0<\min\{1,\varepsilon\}=d$. At $t=1$, it follows from
$c=\max\{0,1-\varepsilon\}<1=d$. For $0<t<1$, both
$0<t$ and $t-\varepsilon<t$ give $c<t$, while
$t<1$ and $t<t+\varepsilon$ give $t<d$.

Every rational number $r\in(c,d)$ satisfies
\[
 0<r<1,
 \qquad t-\varepsilon<r<t+\varepsilon.
\]
It therefore belongs to the enumerated set and satisfies
$\abs{r-t}<\varepsilon$.
By Step~1, there are infinitely many distinct such rational values.
A finite prefix $q_1,\ldots,q_N$ contains only finitely many values,
so it cannot contain them all. Hence for every finite index bound $N$
there is an index $n>N$ such that
\[
 \abs{q_n-t}<\varepsilon.
\]
This reasoning includes the endpoints $t=0$ and $t=1$.

\step{3. Construct a subsequence converging to the chosen target}
Put $n_0=0$. At stage $k\ge1$, apply Step~2 with
$\varepsilon=1/k$ and with the finite index bound $n_{k-1}$.
Choose an index $n_k>n_{k-1}$ satisfying
\[
 \abs{q_{n_k}-t}<\frac1k.
\]
For definiteness, one may take the least eligible index, since a
nonempty set of positive integers has a least element.
The inequalities $n_1<n_2<\cdots$ ensure that $(q_{n_k})$
is a subsequence.

Let $\eta>0$ and choose an integer $K>1/\eta$ by the Archimedean
property. If $k>K$, then
\[
 \abs{q_{n_k}-t}<\frac1k<\eta.
\]
Thus $q_{n_k}\to t$. Since $t\in[0,1]$ was arbitrary, every
point of that interval is a subsequential limit. This is also the
infinite-neighborhood criterion of Theorem~11.2(i).

\step{4. Exclude every other finite or infinite limit}
Every subsequence still satisfies $0<q_{n_k}\le1$.
If it had a finite limit $t<0$, use the positive tolerance
$\eta=-t/2$. Eventually its terms would be less than
$t+\eta=t/2<0$, a contradiction. If it had a finite limit $t>1$,
use $\eta=(t-1)/2$. Eventually its terms would be greater than
$t-\eta=(t+1)/2>1$, again a contradiction. Hence every finite
subsequential limit lies in $[0,1]$.

No subsequence can tend to $+\infty$, because its terms never exceed
$1$ and so cannot eventually exceed the positive threshold $2$.
No subsequence can tend to $-\infty$, because its terms are positive
and so cannot eventually be less than $-1$.
Combining this exclusion with Step~3 proves that the exact set of
subsequential limits is $[0,1]$.

\step{5. Part (b): Take the extreme subsequential limits}
By Theorem~11.8(ii), the limit superior is the supremum of the set
of subsequential limits, and the limit inferior is its infimum.
The interval $[0,1]$ has maximum $1$ and minimum $0$, so
\[
 \limsup_{n\to\infty}q_n=1,
 \qquad \liminf_{n\to\infty}q_n=0.
\]

\answer{\textup{(a)} The set of subsequential limits is $[0,1]$.
\textup{(b)} The limit superior is $1$ and the limit inferior is $0$.
No repeated occurrence of an individual rational value is assumed.}

\exercise{11.6}{A subsequence of a subsequence}
\goal Prove that a subsequence of a subsequence of a given sequence
is itself a subsequence of the original sequence.

\idea Use Definition~11.1 in terms of strictly increasing selection
functions. Compose the two selection functions and verify that the
composition remains strictly increasing.

\step{1. Write the two selections precisely}
Let $s:\N\to\R$ represent the original sequence.
Suppose $t$ is a subsequence of $s$. There is a strictly increasing
function $\sigma:\N\to\N$ with
\[
 t_k=s_{\sigma(k)}\qquad(k\in\N).
\]
Now suppose $u$ is a subsequence of $t$. There is a strictly
increasing function $\tau:\N\to\N$ with
\[
 u_j=t_{\tau(j)}\qquad(j\in\N).
\]
Here strictly increasing refers to the \emph{indices} selected by
the functions; the numerical sequence values themselves need not
increase or be distinct.

\step{2. Form and check the composite selection}
Define $\rho=\sigma\circ\tau$; explicitly,
$\rho(j)=\sigma(\tau(j))$.
Because $\tau$ takes positive integer values and $\sigma$ is defined
on every positive integer, $\rho$ is a function from $\N$ to $\N$.
If $j<\ell$, strict increase of $\tau$ gives
$\tau(j)<\tau(\ell)$. Strict increase of $\sigma$ then gives
\[
 \rho(j)=\sigma(\tau(j))
 <\sigma(\tau(\ell))=\rho(\ell).
\]
Thus $\rho$ is strictly increasing.

\step{3. Identify the selected values}
Substituting the first selection into the second yields, for each $j$,
\[
 u_j=t_{\tau(j)}
     =s_{\sigma(\tau(j))}
     =s_{\rho(j)}.
\]
The values $u_j$ therefore select terms of $s$ in the strictly
increasing index order $\rho(1)<\rho(2)<\cdots$.
Definition~11.1 says exactly that this makes $u$ a subsequence of $s$.

\answer{The composite selection function $\sigma\circ\tau$ proves
that every subsequence of a subsequence is a subsequence of the
original sequence.}

\exercise{11.7}{A subsequence of an enumeration of all rational numbers}
\goal If $(r_n)$ enumerates $\Q$, construct a subsequence
$(r_{n_k})$ tending to $+\infty$.

\idea Every tail of the enumeration remains unbounded above.
Prove that statement with explicit control of a finite prefix,
then select successive terms that exceed their subsequence indices.

\step{1. Show that arbitrarily large rational values occur after any prefix}
Fix a positive integer $N$ and a real threshold $M$.
The finite prefix has a finite maximum
\[
 P=\max\{r_1,\ldots,r_N\}.
\]
By the Archimedean property, choose a positive integer
$p>\max\{M,P,1\}$. Since $p\in\Q$, the enumeration contains $p$:
there is an index $j$ with $r_j=p$.
The inequality $p>P$ prevents $j\le N$, because every term of
that prefix is at most $P$. Thus $j>N$ and $r_j>M$.
This proves that the tail after any finite prefix is unbounded above.
Before any terms are excluded, the same conclusion follows by
choosing an integer $p>M$ in the enumeration.

\step{2. Select terms recursively with increasing indices}
Choose $n_1$ so that $r_{n_1}>1$; this is possible by Step~1.
Suppose $n_1<\cdots<n_{k-1}$ have been chosen.
Apply Step~1 to the prefix ending at $n_{k-1}$ and to the real
threshold
\[
 M_k=\max\{k,r_{n_{k-1}}\}.
\]
It supplies an index $n_k>n_{k-1}$ with
\[
 r_{n_k}>M_k\ge k,
 \qquad r_{n_k}>r_{n_{k-1}}.
\]
The resulting indices are strictly increasing, so they define a
subsequence. The selected values are strictly increasing as well,
although this latter property is not needed just to be a subsequence.

\step{3. Check the definition of divergence to positive infinity}
Let $M>0$ be arbitrary. Choose an integer $K>M$.
If $k>K$, the construction gives
\[
 r_{n_k}>k>K>M.
\]
Thus the selected values eventually exceed every prescribed positive
real threshold. By Definition~9.8, $r_{n_k}\to+\infty$.

\answer{Such a subsequence exists; the construction ensures both
$n_1<n_2<\cdots$ and $r_{n_k}>k$, which proves its limit is $+\infty$.}

\exercise{11.8}{Reflection exchanges the limit inferior and limit superior}
\goal For every real sequence $(s_n)$, prove
\[
 \liminf_{n\to\infty}s_n
 =-\limsup_{n\to\infty}(-s_n)
\]
using Definition~10.6 and the reflection identity of Exercise~5.4.

\idea Reflect each tail before taking its extremum.
If the original sequence is unbounded below, use the infinite-value
conventions in Definition~10.6 directly. Otherwise all relevant tail
bounds are real numbers, and their limits can be compared using
finite convergence or the definition of an infinite limit.

\step{1. Apply reflection to the nonempty tails}
For each positive integer $N$, let
\[
 E_N=\{s_n:n>N\},\qquad
 -E_N=\{-s_n:n>N\}.
\]
Each set is nonempty, since it contains the corresponding term
with index $N+1$. Exercise~5.4 gives the extended equality
\[
 \inf E_N=-\sup(-E_N).
\]
Write $u_N=\inf E_N$ and $v_N=\sup(-E_N)$.
The equality says $u_N=-v_N$, or equivalently $v_N=-u_N$,
with the negation conventions of Section~5.

\step{2. Handle a sequence that is unbounded below}
Suppose $(s_n)$ is unbounded below. Every tail $E_N$ is also
unbounded below. Indeed, if some $E_N$ had a real lower bound $c$,
then
\[
 \min\{c,s_1,\ldots,s_N\}
\]
would be a real lower bound for the entire sequence, a contradiction.
Hence $u_N=-\infty$ for every $N$.

The reflected sequence $(-s_n)$ is unbounded above: for every real
$M$, a term $s_n<-M$ yields $-s_n>M$.
Thus $v_N=+\infty$ for every $N$.
Definition~10.6 explicitly decrees, in these respective cases,
\[
 \liminf s_n=-\infty,
 \qquad \limsup(-s_n)=+\infty.
\]
Consequently
\[
 -\limsup(-s_n)=-(+\infty)=-\infty=\liminf s_n.
\]
This step uses the stated conventions, not a real-sequence limit
law applied to a sequence of infinite symbols.

\step{3. When the original sequence is bounded below, all these tail bounds are real}
Suppose $c\in\R$ is a lower bound for all the $s_n$.
Then every nonempty tail $E_N$ is bounded below by $c$, so
$u_N\in\R$ by Corollary~4.5. The reflected tail is bounded above
by $-c$, so $v_N\in\R$ by completeness. The identity from Step~1
is now an identity of real numbers:
\[
 v_N=-u_N\qquad(N\in\N).
\]
Moreover $E_{N+1}\subseteq E_N$, and restricting the set cannot
decrease its infimum. Thus $u_N\le u_{N+1}$.
The real sequence $(u_N)$ is increasing and is bounded below by $c$.
The monotone-sequence results of Section~10 show that it either
converges to a finite real number or diverges to $+\infty$.
These are the only two cases to consider here.

\step{4. Treat a finite limit of the real tail infima}
If $u_N\to u\in\R$, then Theorem~9.2, applied to the real
sequence $(u_N)$ with scalar $-1$, gives $v_N=-u_N\to-u$.
Definition~10.6 now identifies
\[
 \liminf s_n=u,
 \qquad \limsup(-s_n)=-u.
\]
Negating the second real equality proves
$-\limsup(-s_n)=u=\liminf s_n$.

\step{5. Treat divergence of the real tail infima to positive infinity}
If $u_N\to+\infty$, we show directly that $v_N=-u_N\to-\infty$.
Let $M<0$. The number $-M$ is positive, so the hypothesis supplies
a threshold $N_0$ such that $N>N_0$ implies $u_N>-M$.
Negating gives $v_N=-u_N<M$ for all such $N$.
This is precisely the definition of $v_N\to-\infty$.
Consequently Definition~10.6 gives
\[
 \liminf s_n=+\infty,
 \qquad \limsup(-s_n)=-\infty.
\]
The convention $-(-\infty)=+\infty$ then proves the desired identity.

\answer{In all cases,
$\displaystyle\liminf s_n=-\limsup(-s_n)$.
The infinite cases follow from tail unboundedness and the infinite-limit
definitions, with finite limit laws used only for real tail bounds
having a finite limit.}

\exercise{11.9}{Closed intervals and an impossible set of subsequential limits}
\goal \textup{(a)} Show that $[a,b]$ is closed for real $a\le b$.
\textup{(b)} Decide whether some real sequence can have exactly
$(0,1)$ as its set of subsequential limits.

\idea Use the sequential meaning of closed stated with Theorem~11.9:
the set contains every limit of a convergent sequence of its points.
For part \textup{(b)}, Theorem~11.9 requires the set of subsequential
limits itself to have this property.

\step{1. Part (a): Verify the sequential closedness condition}
Let $(t_k)$ be a sequence with $t_k\in[a,b]$ for every $k$, and
suppose $t_k\to t\in\R$. Thus
\[
 a\le t_k\le b\qquad(k\in\N).
\]
If $t<a$, take $\varepsilon=(a-t)/2>0$.
Eventually $\abs{t_k-t}<\varepsilon$, so
\[
 t_k<t+\varepsilon=\frac{a+t}{2}<a,
\]
contradicting $t_k\ge a$.
If $t>b$, take $\varepsilon=(t-b)/2>0$.
Eventually $\abs{t_k-t}<\varepsilon$, so
\[
 t_k>t-\varepsilon=\frac{b+t}{2}>b,
\]
contradicting $t_k\le b$.
Both alternatives are impossible. Hence $a\le t\le b$, so
$t\in[a,b]$. This proves that $[a,b]$ is closed in the required
sequential sense.

The argument includes $a=b$: in that case all terms equal $a$
and the only possible limit is $a$. If extended limits are also
considered, neither infinite symbol can arise from such a bounded
sequence. For $+\infty$, the positive threshold
$\max\{1,b+1\}>b$ cannot be exceeded; for $-\infty$, the negative
threshold $\min\{-1,a-1\}<a$ cannot be crossed downwards.

\step{2. Part (b): Apply Theorem 11.9 to a sequence of proposed limits}
Suppose, for a contradiction, that a real sequence $(s_n)$ has
exactly $S=(0,1)$ as its set of subsequential limits.
Define
\[
 t_k=\frac1{k+1}\qquad(k\in\N).
\]
Since $k+1\ge2$, we have $0<t_k\le1/2<1$. Thus $t_k\in S$
for every $k$.
Also $t_k\to0$: for any $\varepsilon>0$, choose an integer
$K>1/\varepsilon$; then $k>K$ gives
\[
 \abs{t_k-0}=\frac1{k+1}<\frac1k<\varepsilon.
\]
Theorem~11.9 applies to this sequence of finite members of $S$
and forces its limit $0$ to belong to $S$.
But $0\notin(0,1)$, a contradiction.

\answer{\textup{(a)} Every closed interval $[a,b]$, including a
one-point interval, is closed in the sequential sense.
\textup{(b)} No sequence has $(0,1)$ as exactly its set of
subsequential limits.}

\exercise{11.10}{The alternating diagonal traversal in Figure 11.2}
\goal For the sequence listed by the arrows of Figure~11.2, find
\textup{(a)} its exact set of subsequential limits and
\textup{(b)} its limit superior and limit inferior.

\idea The figure has the same row repeated infinitely many times.
Each reciprocal value therefore occurs at arbitrarily large indices.
The top row also supplies a subsequence of values tending to $0$.
To prove completeness of the answer, isolate any other positive
proposed limit from all reciprocal values.

\step{1. Describe the array and its order precisely}
Number rows and columns by $i,j\in\N$. The entry in position
$(i,j)$ is $1/j$, so the array is
\[
 \begin{pmatrix}
 1&\frac12&\frac13&\frac14&\frac15&\cdots\\
 1&\frac12&\frac13&\frac14&\frac15&\cdots\\
 1&\frac12&\frac13&\frac14&\frac15&\cdots\\
 1&\frac12&\frac13&\frac14&\frac15&\cdots\\
 \vdots&\vdots&\vdots&\vdots&\vdots&\ddots
 \end{pmatrix}.
\]
The traversal groups the entries by diagonals
$i+j=m+1$, with $m=1,2,3,\ldots$.
Diagonal $m$ has exactly $m$ entries.
For odd $m$, the traversal lists columns $1,2,\ldots,m$;
for even $m$, it lists columns $m,m-1,\ldots,1$.
Thus the successive diagonal blocks begin
\[
 \underbrace{1}_{m=1},\quad
 \underbrace{\tfrac12,1}_{m=2},\quad
 \underbrace{1,\tfrac12,\tfrac13}_{m=3},\quad
 \underbrace{\tfrac14,\tfrac13,\tfrac12,1}_{m=4},\quad\ldots.
\]
This order agrees with the alternating directions of the arrows.

\step{2. Check explicitly that choosing later diagonals gives later indices}
Before diagonal $m$ there are
\[
 T_{m-1}=1+2+\cdots+(m-1)=\frac{m(m-1)}2
\]
listed entries, with $T_0=0$.
For $m\ge j$, the occurrence of $1/j$ in diagonal $m$ has index
\[
 I(m,j)=
 \begin{cases}
 \dfrac{m(m-1)}2+j,&m\text{ odd},\\[4pt]
 \dfrac{m(m-1)}2+(m-j+1),&m\text{ even}.
 \end{cases}
\]
Indeed, its position within that diagonal is $j$ when columns are
in increasing order, and $m-j+1$ when they are in decreasing order.
In both cases
\[
 T_{m-1}+1\le I(m,j)\le T_m,
 \qquad T_m=\frac{m(m+1)}2.
\]
Every index of diagonal $m+1$ is at least $T_m+1$, which is
strictly larger than every index of diagonal $m$.
Consequently selections made from successive diagonals always
have strictly increasing indices.

\step{3. Part (a): Obtain every reciprocal as a subsequential limit}
Fix $j\in\N$. Every diagonal $m\ge j$ contains column $j$
once, at the matrix position $(m+1-j,j)$.
Choose the indices
\[
 n_k=I(j+k-1,j)\qquad(k\in\N).
\]
They are strictly increasing by Step~2, and the corresponding
sequence values satisfy $s_{n_k}=1/j$ for every $k$.
This constant subsequence converges to $1/j$.
Thus every member of $\{1/j:j\in\N\}$ is a subsequential limit.

\step{4. Part (a): Obtain zero from the top row}
The top-row entry in column $k$ lies in diagonal $k$ and has
value $1/k$. Select it at index
\[
 n_k=I(k,k)=
 \begin{cases}
 \dfrac{k(k+1)}2,&k\text{ odd},\\[4pt]
 \dfrac{k(k-1)}2+1,&k\text{ even}.
 \end{cases}
\]
These indices are strictly increasing because they come from
successive diagonals. For example, they begin $1,2,6,7,15,\ldots$.
The selected values are exactly $s_{n_k}=1/k$.
Given $\varepsilon>0$, an integer $K>1/\varepsilon$ ensures
$k>K$ implies $\abs{1/k-0}<\varepsilon$.
Therefore $s_{n_k}\to0$, and $0$ is also a subsequential limit.

\step{5. Part (a): Exclude every other proposed limit}
Every term of the full sequence belongs to
$D=\{1/j:j\in\N\}\subseteq(0,1]$.
Hence any finite subsequential limit is in $[0,1]$, by the
same eventual-bound argument used in Exercise~11.5.
The bounds also rule out both infinite subsequential limits.

Suppose $t>0$ and $t\notin D$. We show that a fixed positive
neighborhood of $t$ contains no term of the sequence.
Choose an integer $J>\max\{2,2/t\}$.
For every $j\ge J$, we have
\[
 \frac1j\le\frac1J<\frac t2,
 \qquad \abs{\frac1j-t}=t-\frac1j>\frac t2.
\]
For the finitely many indices $1\le j<J$, none of the values
$1/j$ equals $t$. Their distances from $t$ are strictly positive,
so the finite minimum
\[
 d=\min_{1\le j<J}\abs{\frac1j-t}
\]
exists and is positive. Define
$\delta=\min\{t/4,d/2\}>0$.
If $j\ge J$, then $\abs{1/j-t}>t/2>\delta$.
If $j<J$, then $\abs{1/j-t}\ge d>\delta$.
Thus every term of the original sequence satisfies
$\abs{s_n-t}>\delta$.
No subsequence can converge to $t$, since convergence would require
its terms eventually to have distance less than $\delta$ from $t$.

Every possible finite limit is either $0$, a member of $D$, or
a positive number not in $D$. Steps~3--5 account for all these
possibilities, and the infinite cases were excluded. Therefore
\[
 S=\{0\}\cup\left\{\frac1j:j\in\N\right\}
\]
is exactly the set of subsequential limits.

\step{6. Part (b): Use the extreme members of the computed set}
The set $S$ has minimum $0$ and maximum $1$.
Theorem~11.8(ii) therefore gives
\[
 \limsup_{n\to\infty}s_n=\sup S=1,
 \qquad \liminf_{n\to\infty}s_n=\inf S=0.
\]

\answer{\textup{(a)} The exact set is
$\displaystyle S=\{0\}\cup\{1/j:j\in\N\}$.
\textup{(b)} The limit superior is $1$ and the limit inferior is $0$.}

\exercise{11.11}{An increasing sequence approaching a supremum from the set}
\goal For a nonempty bounded set $S\subseteq\R$, construct an
increasing sequence $(s_n)$ with every $s_n\in S$ and
$s_n\to\sup S$.

\idea If the supremum belongs to the set, use the constant sequence
at that point. Otherwise, at each stage choose a point of $S$ above
both the previous value and a closer lower approximation to the
supremum.

\step{1. State the necessary nonemptiness condition and the meaning of increasing}
The exercise's statement needs the implicit assumption $S\ne\varnothing$.
For the empty set there is no sequence of points in $S$, and its
supremum is not defined by the real-supremum conventions being used.
We therefore assume that $S$ is nonempty and bounded.
In fact, the construction needs only that it be bounded above.
Completeness gives the finite real number
\[
 M=\sup S.
\]
As in Section~10 and in the note to this exercise, increasing means
nondecreasing: $s_n\le s_{n+1}$. A constant sequence is allowed.
When $M\notin S$, the construction below will give strict increase.

\step{2. The case in which the supremum belongs to the set}
If $M\in S$, define $s_n=M$ for every $n\in\N$.
Then all terms belong to $S$ and $s_n=s_{n+1}$, so the sequence
is increasing in the stated sense. For every $\varepsilon>0$,
\[
 \abs{s_n-M}=0<\varepsilon
\]
at every index. Hence $s_n\to M$.

\step{3. Begin the construction when $M$ is not a member}
Suppose $M\notin S$. Since $M$ is an upper bound, every $s\in S$
satisfies $s\le M$; the nonmembership of $M$ strengthens this to
$s<M$ for every $s\in S$.
The number $M-1<M$ cannot be an upper bound for $S$, because
$M$ is its least upper bound. Therefore some $s_1\in S$ satisfies
\[
 M-1<s_1<M.
\]

\step{4. Continue the construction with an explicitly admissible threshold}
Suppose $s_{n-1}\in S$ has been chosen, where $n\ge2$.
Both $s_{n-1}<M$ and $M-1/n<M$ hold. Therefore
\[
 c_n=\max\left\{s_{n-1},M-\frac1n\right\}<M.
\]
As $c_n<M=\sup S$, the number $c_n$ is not an upper bound for $S$.
Choose $s_n\in S$ with $s_n>c_n$.
Every point of $S$ is less than $M$, so this choice satisfies
\[
 s_{n-1}<s_n<M,
 \qquad M-\frac1n<s_n<M.
\]
The procedure can be repeated at every positive integer stage.
It gives a sequence in $S$ that is strictly increasing.
The displayed approximation bound also holds at $n=1$ by Step~3.

\step{5. Verify convergence to the supremum}
For every $n$, the construction gives
\[
 0<M-s_n<\frac1n.
\]
Let $\varepsilon>0$. Choose an integer $N>1/\varepsilon$ by the
Archimedean property. If $n>N$, then
\[
 \abs{s_n-M}=M-s_n<\frac1n<\varepsilon.
\]
Thus $s_n\to M$ by the definition of convergence.

\answer{With the necessary hypothesis $S\ne\varnothing$, an increasing
sequence in $S$ converging to $\sup S$ always exists. It may be taken
constant when $\sup S\in S$, and strictly increasing when
$\sup S\notin S$.}
